\chapter{Problem Formulation}

\section{Geometry Model}
The cooling tower is modeled as a stack of frustums.
For segment $i$ with lower and upper radii $r_{i-1}, r_i$ and height $h_i$:
\begin{align}
A_i &= \pi (r_{i-1} + r_i)\sqrt{(r_i-r_{i-1})^2 + h_i^2},\\
V_i &= \frac{\pi h_i}{3}\left(r_{i-1}^2 + r_{i-1}r_i + r_i^2\right).
\end{align}
Total shell area and volume are:
\begin{align}
\towerArea &= \sum_i A_i,\quad
\towerVolume = \sum_i V_i.
\end{align}

\section{Optimization Objective}
The scalar objective minimized by all solvers is:
\begin{equation}
\obj(x) = \towerArea(x) + P_{\mathrm{volume}} + P_{\mathrm{height}} + P_{\mathrm{shape}} + P_{\mathrm{smooth}} + P_{\mathrm{bounds}} + P_{\mathrm{ratio}},
\end{equation}
with terms activated per scenario.

\section{Feasibility Criteria}
A run is feasible if all relevant conditions are satisfied:
\begin{itemize}[leftmargin=1.5em]
\item Relative volume error $\le 10^{-3}$,
\item Relative total-height error $\le 10^{-3}$ when heights are decision variables,
\item Monotonic hyperboloid violation $\le 10^{-6}$ when radii are decision variables.
\end{itemize}
